% !PRINTED 492
\weberpage{492}
%<<CONT
mit $\psi$ nicht identisch. Man kann daher (nach \S.\,143, 1.) dem $t$ einen solchen rationalen Werth geben, dass $\psi$ von allen nach Art von $\psi_1$ gebildeten Functionen verschieden ist.

Selbstverständlich gelten diese Betrachtungen auch, wenn an Stelle der $\alpha$ unabhängige Variable treten.

Eine Function, die alle Permutationen einer Gruppe $P$ gestattet, während sie sich bei allen nicht zu $P$ gehörigen Permutationen ändert, heisst eine zur \gesperrt{Gruppe $P$ gehörige Function}.

Zu der Gruppe $P_0$, die \gesperrt{alle} Permutationen von $m$ Ziffern umfasst, gehören die symmetrischen Functionen. Man nennt diese Gruppe daher auch die \gesperrt{symmetrische Gruppe}.

Eine Function $\varrho$, die zu der aus der einzigen identischen Permutation bestehenden Gruppe gehört, kann zu einer eindeutigen Bezeichnungsweise der Permutationen dienen; denn es giebt nur eine Permutation $\pi'$, durch die $\varrho$ in $\varrho'$ übergeht. Wir können also für diese Permutation das Symbol
\[
\pi' = (\varrho, \varrho')
\]
benutzen, so dass durch
\[
(\varrho, \varrho), \quad (\varrho, \varrho'), \quad (\varrho, \varrho'') \dots
\]
alle Permutationen der $m$ Ziffern eindeutig bezeichnet sind.

\begin{center}
\S.\,153.

\gesperrt{Zerlegung von Permutationen in Transpositionen\\
und in Cyklen.}
\end{center}

Wir haben schon im zweiten Abschnitt bei Gelegenheit der Determinanten die Zusammensetzung beliebiger Permutationen aus einer Reihe von Transpositionen, d. h. Vertauschung nur zweier Ziffern, erwähnt. Solche Transpositionen bezeichnen wir durch Nebeneinanderstellen der betreffenden Ziffern, z. B. durch $(0,1)$ die Vertauschung der Ziffern 0 und 1.

Eine Transposition, zweimal wiederholt, führt zur Identität, so dass jede Transposition sich selbst entgegengesetzt ist. Ist nun
\[
\pi = \begin{pmatrix} 0, & 1, & 2 \dots m-1 \\ a_0, & a_1, & a_2 \dots a_{m-1} \end{pmatrix}
\]
eine beliebige Permutation von $m$ Ziffern, so ist
\[
\pi(m-1, a_{m-1}) = \pi_1
\]
eine Permutation, die $m-1$ ungeändert lässt, die also durch
%CONT>>
